\documentclass[10pt,a4paper]{amsart}
\usepackage[margin=19mm]{geometry}
\usepackage{amsmath,amssymb,mathtools}
\usepackage[scr=rsfs,cal=euler]{mathalfa}
\usepackage{xfrac}
\usepackage[unicode,hidelinks,bookmarks=false]{hyperref}
\newcommand{\ie}{i.e.\ }
\newcommand{\cf}{cf.\ }
\newcommand{\resp}{respectively\ }
\newcommand{\Subsection}{\subsection}
\providecommand{\nobreakdash}{\nobreak\hbox{-}}
\setlength{\emergencystretch}{2em}
\multlinegap0pt
\usepackage{amssymb}
\usepackage{mathtools}
\newtheorem{thm}{Theorem}
\title{Extensions of several famous combinatorial identities via hypergeometric functions}
\author{Arjun Kumar Rathie, Feng Qi, Dongkyu Lim*}
\date{Source: arXiv:2608.05192v2 (CC BY 4.0)}
\begin{document}
\maketitle

\noindent\emph{Source and licence.} Extensions of several famous combinatorial identities via hypergeometric functions. Version arXiv:2608.05192v2 (CC BY 4.0), distributed by arXiv under the Creative Commons Attribution 4.0 International licence, http://creativecommons.org/licenses/by/4.0/, as recorded in the arXiv licence metadata for this exact version and shown on the abstract page https://arxiv.org/abs/2608.05192v2. Full licence notice: \texttt{licenses/arxiv-2608.05192-CC-BY-4.0.txt}.\par\smallskip

\noindent\textit{Editorial scope.} Counted unit: the article's thm 5.1thm (source0.tex lines 794--807). The article's complete original proof of it is delivered with it (source0.tex lines 809--833, one \texttt{proof} environment of 25 lines). No mathematics is written, repaired or re-derived: the statement and the proof are the article's own lines, in the article's own notation, and no retained line is substituted or re-flowed.  Retained uncounted as the source-local context the proof needs: lines 241--248; lines 327--335.  Nothing was omitted from any retained span, so the package carries no omission record.  No retained line contains a citation, so the wrapper carries no bibliography and no bibliography entry.  No retained line contains a figure environment, an image-inclusion command or drawing code, so no image is delivered and the package declares no asset dependency.  Editorial formatting only, in addition: an AMS \texttt{amsart} preamble that restates the article's own declarations and macros used by the retained lines, namely the environments \texttt{thm} on separate counters, together with the standard packages those lines need.  The retained environments print in this document as Theorem 1 in order of appearance, whereas the article prints the same items with the numbering of its own full document.  The complete native source of the article as distributed in the arXiv e-print for this exact version is bundled unchanged as \texttt{sources/206909/source0.tex}.
\begin{align}\label{15}
{\,}_3F_2\begin{bmatrix}\begin{gathered}-2n, a, d+1\\ 2a+1, d \end{gathered};2\end{bmatrix}
&=\frac{\bigl(\frac{1}{2}\bigr)_n}{\bigl(a+\frac{1}{2}\bigr)_n}
\intertext{and}
\label{16}
{\,}_3F_2\begin{bmatrix}\begin{gathered}-2n-1, a, d+1\\ 2a+1, d \end{gathered};2\end{bmatrix}
&=\frac{\bigl(1-\frac{2a}{d}\bigr)\bigl(\frac{3}{2}\bigr)_n}{\bigl(2a+1\bigr)\bigl(a+\frac{3}{2}\bigr)_n}
\end{align}

\begin{align}\label{25}
\binom{2n}{n}&=2^{2n}\frac{\bigl(\frac{1}{2}\bigr)_n}{(1)_n}, & \binom{n}{k}&=\frac{(-1)^k(-n)_k}{(1)_k},\\
\label{27}
\binom{k+n}{k}&=\frac{\bigl(n+1\bigr)_k}{(1)_k}, & \binom{n+1}{k+1}&=(-1)^k(n+1)\frac{\bigl(-n\bigr)_k}{(2)_k},\\
\label{29}
\binom{k+d}{k+1}&=\frac{d(d+1)_k}{(2)_k}, & \binom{n}{2k}&=\frac{\bigl(-\frac{n}2\bigr)_k\bigl(-\frac{n}2+\frac12\bigr)_k}{\bigl(\frac12\bigr)_k(1)_k},\\
\label{26}
\binom{k+d}{k}&=\frac{(d+1)_k}{(1)_k}, & C_n&=2^{2n}\frac{\bigl(\frac{1}{2}\bigr)_n}{(2)_n}.
\end{align}

\begin{thm}\label{5.1thm}
For $d \neq 0,-1,-2,\dotsc$ and $n,\nu \in \mathbb{N}_0$, we have
\begin{equation}\label{(5.1)}
\sum_{k=0}^{n}\frac{(-1)^k}{2^k(k+1)}
\binom{n}{k}\binom{2k}{k}
\frac{\binom{k+d}{k}}{\binom{k+d-1}{k}}
=
\begin{dcases}
\frac{\bigl(\frac12\bigr)_\nu}{(1)_\nu}, & n=2\nu; \\
\frac12\biggl(1-\frac{1}{d}\biggr)
\frac{\bigl(\frac32\bigr)_\nu}{(2)_\nu}, & n=2\nu+1.
\end{dcases}
\end{equation}
\end{thm}

\begin{proof}
Denote the left-hand side of~\eqref{(5.1)} by $S_1$. Converting all binomial coefficients involved in $S_1$ to the Pochhammer symbols by employing the appropriate elementary identities~\eqref{25} to~\eqref{26} results in
\[
S_1
= \sum_{k=0}^{n}
\frac{(-n)_k\bigl(\frac12\bigr)_k(d+1)_k}{(2)_k(d)_k}\frac{2^k}{k!}
= {}_3F_2
\begin{bmatrix}
\begin{gathered}-n,\tfrac12,d+1 \\
2,d\end{gathered}; 2
\end{bmatrix}.
\]
Taking $a=\frac12$ in~\eqref{15} and~\eqref{16} immediately leads to
\[
{}_3F_2\begin{bmatrix}\begin{gathered}-n, \tfrac{1}{2}, d+1\\
2,d\end{gathered};2\end{bmatrix}
=
\begin{dcases}
\frac{\bigl(\frac12\bigr)_\nu}{(1)_\nu}, & n=2\nu; \\
\frac12\biggl(1-\frac{1}{d}\biggr)
\frac{\bigl(\frac32\bigr)_\nu}{(2)_\nu}, & n=2\nu+1.
\end{dcases}
\]
Accordingly, the identity~\eqref{(5.1)} is proved. Theorem~\ref{5.1thm} is thus proved.
\end{proof}

\end{document}
